\section*{अध्याय \ref{ch:b2:blood-pressure} --- रक्तचाप का नियमन और व्यायाम}

\begin{solution}{exo:b2:blood-pressure:1}
$\bar P - P_{\text{शिरा}} = Q\,R = f\,V_s\,R$। हृद्-दर (यानी
\omterm{prop:b2:heart:conduction}{साइनु-अलिंद गाँठ}, वेगस तथा अनुकंपी तंत्रिकाओं के नीचे);
\omterm{prop:b2:heart:cycle}{आघात-आयतन} (यानी वह \omterm{def:b2:heart:anatomy}{निलय}: उसका भरना, जो \omterm{def:b2:blood-circulation:vessels}{शिराएँ} तय करती हैं, और उसकी
सुकुंचनशीलता, जो अनुकंपी तंत्रिकाएँ तय करती हैं); प्रतिरोध (यानी
\omterm{def:b2:blood-circulation:vessels}{धमनिका} की चिकनी पेशी, अनुकंपी तंत्रिकाओं, हॉर्मोनों तथा
स्थानीय उपापचयजों के नीचे)।
\end{solution}

\begin{solution}{exo:b2:blood-pressure:2}
संवेदक: ग्रीवा ज्यावों तथा महाधमनी-चाप में खिंचाव-ग्राही।
केंद्र: मेडुला के हृद्-वाहिका केंद्र। प्रभावक:
वेगस (दर) और अनुकंपी तंत्रिकाएँ (दर, सुकुंचनशीलता,
\omterm{def:b2:blood-circulation:vessels}{धमनिका} तथा \omterm{def:b2:blood-circulation:vessels}{शिरा} की तान)। ऋण प्रतिपुष्टि, लगभग एक सेकंड की
देरी। उन संवेदकों के बिना दिन भर का माध्य दाब
सामान्य रहता है पर मिनट-दर-मिनट दाब हर मुद्रा तथा भाव के साथ
जंगली ढंग से झूलता है।
\end{solution}

\begin{solution}{exo:b2:blood-pressure:3}
वृक्क (धमनिका-कोशिकाएँ) रेनिन तब छोड़ता है जब उसका सींचा जाना या
सोडियम-पहुँच गिरे या उसकी अनुकंपी तंत्रिकाएँ दगें; रेनिन
ऐंजियोटेंसिनोजन (यकृत) को काटकर ऐंजियोटेंसिन I बनाता है; फेफड़े का एंजाइम उसे
ऐंजियोटेंसिन II में बदलता है, जो \omterm{def:b2:blood-circulation:vessels}{धमनिकाएँ} सिकोड़ता है और
अधिवृक्क वल्कुट को ऐल्डोस्टेरॉन स्रवित करने को उकसाता है; ऐल्डोस्टेरॉन वृक्क से
सोडियम तथा जल रुकवाता है, और रक्त-आयतन चढ़ा देता है।
\end{solution}

\begin{solution}{exo:b2:blood-pressure:4}
उस पेशी की अपनी \omterm{def:b2:blood-circulation:vessels}{धमनिकाओं} का उसके उपापचयजों से उपापचयी
प्रसार; अनुकंपी तंत्रिकाओं तथा पेशी-पंप से चलाया गया
पाँच गुना \omterm{thm:b2:heart:starling}{हृद्-निर्गम}; और आँत, वृक्कों तथा विश्राम करती पेशी की
\omterm{def:b2:blood-circulation:vessels}{धमनिकाओं} का अनुकंपी संकोचन, जो प्रवाह छोड़
देती हैं।
\end{solution}

\begin{solution}{exo:b2:blood-pressure:5}
$83\times 1.2 = \qty{100}{mmHg}$; $333\times 0.35 = \qty{117}{mmHg}$।
प्रतिरोध $1.2/0.35 = 3.4$ से नीचे: त्रिज्या $3.4^{1/4} =
1.36$ से ऊपर।
\end{solution}

\begin{solution}{exo:b2:blood-pressure:6}
$30/(1 + G)$: 10, 6 तथा \qty{3.3}{mmHg}। आधी लब्धि के साथ वह
बचा हुआ पात दुगुना हो जाता है: वह खड़े होने पर चकराती है, धूसर देखती है, और
मूर्छित हो सकती है।
\end{solution}

\begin{solution}{exo:b2:blood-pressure:7}
$\rho g h$: $1060\times 9.81\times 0.4 = \qty{4160}{Pa} =
\qty{31}{mmHg}$; $1060\times 9.81\times 1.2 = \qty{12500}{Pa} =
\qty{94}{mmHg}$। मस्तिष्क $95 - 31 = \qty{64}{mmHg}$; टख़ना $95 + 94 =
\qty{189}{mmHg}$। किसी जिराफ़ का \qty{2.5}{m} पर सिर
\qty{195}{mmHg} लेता है: उसे \omterm{def:b2:heart:anatomy}{हृदय} पर कोई \qty{250}{mmHg} का माध्य दाब
चाहिए, और टाँगों में मोटी वाहिका-भित्तियाँ तथा कसी त्वचा।
\end{solution}

\begin{solution}{exo:b2:blood-pressure:8}
$Q = 280/(190 - 130) = \qty{4.7}{L/min}$। खिलाड़ी: $5000/(200 - 30) =
\qty{29}{L/min}$; \omterm{prop:b2:heart:cycle}{आघात-आयतन} $\num{29400}/190 = \qty{155}{mL}$।
\end{solution}

\begin{solution}{exo:b2:blood-pressure:9}
$35/5 = \qty{7}{mmHg}$ नीचे: यानी लगभग \qty{88}{mmHg}। उस
आयतन की बहाली: अंतरालीय तरल का \omterm{def:b2:blood-circulation:vessels}{केशिकाओं} में फिर सोखा जाना
(मिनटों से घंटे भर में); \omterm{prop:b2:blood-pressure:raas}{प्रतिमूत्रल हॉर्मोन} तथा
ऐल्डोस्टेरॉन के नीचे घटा मूत्र (घंटों से दिनों में); प्यास तथा पीना (घंटों में); मज्जा से
लाल कोशिकाएँ (सप्ताहों में)। त्वचा पीली तथा ठंडी है क्योंकि उस
प्रतिवर्त ने मस्तिष्क तथा \omterm{def:b2:heart:anatomy}{हृदय} के लिए वह दाब बचाने को उसकी \omterm{def:b2:blood-circulation:vessels}{धमनिकाएँ} बंद
कर दी हैं।
\end{solution}

\begin{solution}{exo:b2:blood-pressure:10}
$R = 0.5$ पर पेशी: $95/0.5 = \qty{190}{mL/s}$, \qty{11.4}{L/min};
बाक़ी अंगों के \qty{4}{L/min} के साथ \omterm{def:b2:heart:anatomy}{हृदय} को
\qty{15.4}{L/min} चाहिए होता। $Q$ \qty{5}{L/min} पर नियत हो, तो कुल
प्रतिरोध 1.14 से गिरकर $1/(1/1.14 - 1/5 + 1/0.5) =
\qty{0.37}{mmHg.s/mL}$ रह जाता है और दाब \qty{31}{mmHg} --- और
मस्तिष्क मूर्छित हो जाता। देह निर्गम पहली स्थिति की ओर चढ़ाती है
और बाक़ी शय्याएँ सिकोड़ती है, ताकि वह दाब थमा रहे और वह
पेशी अपना प्रवाह पा ले।
\end{solution}

\begin{solution}{exo:b2:blood-pressure:11}
(a) कम सींचे वृक्क ने कोई नीचा दाब पढ़ा, रेनिन छोड़ा,
और ऐंजियोटेंसिन तथा ऐल्डोस्टेरॉन ने वह प्रतिरोध तथा वह
आयतन चढ़ा दिया। (b) \omterm{def:b2:blood-pressure:baroreflex}{दाब-प्रतिवर्त} परिवर्तनों को कुंद करता है पर जो भी
स्तर बना रहे उसी पर पुनः बैठ जाता है, और वृक्क के आयतन रोकने को बदल नहीं सकता।
(c) रेनिन का स्रोत, और ऊँचा दाब माँगता अंग,
जा चुका था। (d) ऐंजियोटेंसिन II में रूपांतरण रोक देता: नीचा
प्रतिरोध, कम ऐल्डोस्टेरॉन, नीचा दाब।
\end{solution}

\begin{solution}{exo:b2:blood-pressure:12}
\omterm{def:b2:blood-pressure:baroreflex}{दाब-प्रतिवर्त} एक सेकंड में लगभग चार की लब्धि और ऐसे नियत
बिंदु के साथ काम करता है जो प्रचलित दाब की ओर सरक जाता है: वह उतार-चढ़ाव कुंद करता है
और कोई स्तर थाम नहीं सकता। वृक्क घंटों तथा दिनों भर काम करता है, और
प्रभाव में अनंत लब्धि के साथ --- वह तब तक निकालता या रोकता रहता है जब तक
वह दाब न आ जाए जिस पर वह बैठा है --- और उसका नियत बिंदु उसकी अपनी
शरीरक्रिया से नियत है: वह उस स्तर को थामता है। इसलिए कोई चिरकालिक चढ़ाई
का अर्थ है कि वृक्क का नियत बिंदु सरक चुका है, और उसे केवल वे दवाएँ नीचे लाती हैं जो
वृक्क के पाश पर या उसके माँगे प्रतिरोध पर काम करें।
\end{solution}

\begin{solution}{pb:b2:blood-pressure:1}
\textbf{1.} $1/R = 1/7.6 + 1/22.8 + 1/4.1 + 1/5.2 + 1/5.7 + 1/19 +
1/28.5 = 0.875$: $R = \qty{1.14}{mmHg.s/mL}$; $\bar P/Q = 95/83.3 =
1.14$।
\textbf{2.} $\bar P/R_i$: मस्तिष्क 0.75, \omterm{def:b2:heart:anatomy}{हृदय} 0.25, आँत तथा यकृत
1.39, वृक्क 1.10, पेशी 1.0, त्वचा 0.30, अन्य \qty{0.20}{L/min};
अंश 15, 5, 28, 22, 20, 6 तथा \qty{4}{\%}।
\textbf{3.} $5000/70 = \qty{71}{mL}$; प्रति मिनट $5\times 50 = \qty{250}{mL}$
ऑक्सीजन।
\textbf{4.} मस्तिष्क $0.75/1.4 = \qty{0.54}{L/min/kg}$; वृक्क $1.1/0.3
= \qty{3.7}{L/min/kg}$, यानी मस्तिष्क से सात गुना और देह के
औसत से पचास गुना: वृक्क छानने के लिए सींचे जाते हैं, पोसने के लिए नहीं।
\textbf{5.} जब एक ही दाब सारे अंगों में साझा हो, तब हर एक वही
प्रवाह खींचता है जो उसका अपना प्रतिरोध तय करे; प्रवाहों को केंद्र से
नियंत्रित करना किसी भी ऐसे अंग को भूखा रखता जिसकी माँग बदली हो। दाब ही
साझी मुद्रा है।
\textbf{6.} \omterm{def:b2:heart:anatomy}{हृदय}: $250\times 0.2\times 0.7 = \qty{35}{mL/min}$;
पेशी: $1000\times 0.2\times 0.25 = \qty{50}{mL/min}$। \omterm{def:b2:heart:anatomy}{हृदय}
जो गुज़रता है उसका अधिकांश पहले ही निकाल लेता है, इसलिए वह ऑक्सीजन केवल
अधिक प्रवाह से पा सकता है, और इसीलिए जैसे ही वह अधिक काम करता है उसकी \omterm{def:b2:blood-circulation:vessels}{धमनिकाएँ} फैल जाती हैं।
\textbf{7.} $Q$ \qty{30}{\%} नीचे, नियत $R$ पर: $\bar P$
\qty{30}{\%} नीचे, यानी \qty{28.5}{mmHg}, गिरकर \qty{66.5}{mmHg}।
\textbf{8.} $28.5/5 = \qty{5.7}{mmHg}$: यानी लगभग \qty{89}{mmHg}।
\textbf{9.} $Q = 85\times 0.7\times 71 = \qty{4.2}{L/min} =
\qty{70}{mL/s}$; $R = 89.3/70 = \qty{1.27}{mmHg.s/mL}$, यानी
\qty{11}{\%} ऊपर।
\textbf{10.} $1060\times 9.81\times 0.4 = \qty{4160}{Pa} =
\qty{31}{mmHg}$; मस्तिष्क उस प्रतिवर्त से पहले $66.5 - 31 = \qty{35}{mmHg}$,
और बाद $89 - 31 = \qty{58}{mmHg}$।
\textbf{11.} $28.5/2 = \qty{14}{mmHg}$: \omterm{def:b2:heart:anatomy}{हृदय} पर \qty{81}{mmHg},
मस्तिष्क पर \qty{50}{mmHg} --- यानी चक्कर, धूसर दृष्टि, और कुर्सी से
हर बार उठने पर मूर्छा के क़रीब।
\textbf{12.} सीधा लेटना \qty{31}{mmHg} की द्रवस्थैतिक हानि हटा देता है
और इकट्ठे \omterm{def:b2:blood-circulation:blood}{रक्त} को \omterm{def:b2:heart:anatomy}{हृदय} तक वापस बहा देता है; \omterm{prop:b2:heart:cycle}{आघात-आयतन} और
मस्तिष्क-दाब कुछ ही धड़कनों में लौट आते हैं।
\textbf{13.} $45/5 = \qty{9}{mmHg}$: \qty{86}{mmHg}।
\textbf{14.} $1/R = 0.875 - 1/19 - 1/4.1 + 1/38 + 1/8.2 = 0.727$: $R =
\qty{1.38}{mmHg.s/mL}$। मस्तिष्क $86/7.6 = \qty{11.3}{mL/s} =
\qty{0.68}{L/min}$; त्वचा $86/38 = \qty{2.3}{mL/s} = \qty{0.14}{L/min}$;
आँत $86/8.2 = \qty{10.5}{mL/s} = \qty{0.63}{L/min}$।
\textbf{15.} कम \omterm{def:b2:blood-circulation:blood}{रक्त} तथा कसी \omterm{def:b2:blood-circulation:vessels}{धमनिकाओं} के साथ
केशिका-दाब पूरी \omterm{def:b2:blood-circulation:vessels}{केशिका} भर \omterm{thm:b2:blood-circulation:oncotic}{परासरणी दाब} से नीचे गिर जाता है,
इसलिए तरल छनने के बजाय फिर सोखा जाता है; \omterm{def:b2:blood-circulation:blood}{प्लाज़्मा} तनु होता है और
\omterm{def:b2:blood-circulation:blood}{हीमैटोक्रिट} गिरता है।
\textbf{16.} ऐल्डोस्टेरॉन, रेनिन तथा ऐंजियोटेंसिन से होकर, जिसे वृक्क का
नीचा सींचा जाना तथा अनुकंपी चालन छेड़ते हैं; \omterm{prop:b2:blood-pressure:raas}{प्रतिमूत्रल हॉर्मोन},
जिसे आयतन की गिरावट तथा प्लाज़्मा-सांद्रता की चढ़ाई
छेड़ती है।
\textbf{17.} $5\times 10^{12}$ कोशिकाएँ $2.5\times 10^{6}$ प्रति सेकंड पर:
$2\times 10^{6}$ s, यानी लगभग तीन सप्ताह।
\textbf{18.} लगभग \qty{60}{mmHg} से नीचे वह प्रतिवर्त-वक्र सपाट है:
वाहिकाएँ जितनी कस सकती थीं और \omterm{def:b2:heart:anatomy}{हृदय} जितना तेज़ हो सकता था, हो चुका, और लब्धि
शून्य है, इसलिए हर आगे की हानि उस दाब को अपनी पूरी मात्रा भर
नीचा करती है; और तब कम सींचे ऊतक अम्ल छोड़कर फैल जाते हैं, और
वह दाब ढह जाता है।
\textbf{19.} पेशी $110/0.33 = \qty{333}{mL/s} = \qty{20}{L/min}$;
त्वचा $110/5 = \qty{22}{mL/s} = \qty{1.3}{L/min}$।
\textbf{20.} मस्तिष्क-प्रतिरोध $8.8$ (110 पर प्रवाह 0.75), \omterm{def:b2:heart:anatomy}{हृदय} $6.6$
(प्रवाह 1.0): $1/R = 1/8.8 + 1/6.6 + 1/8.2 + 1/10.4 + 1/0.33 + 1/5 +
1/28.5 = 3.75$: $R = \qty{0.27}{mmHg.s/mL}$; $Q = 110/0.27 =
\qty{410}{mL/s} = \qty{25}{L/min}$।
\textbf{21.} $\num{24700}/180 = \qty{137}{mL}$।
\textbf{22.} $24.7\times(200 - 40) = \qty{3950}{mL/min}$, यानी
विश्राम के \qty{250}{mL/min} से सोलह गुना।
\textbf{23.} प्रवाह $\times 4.9$, निष्कर्षण $160/50 = \times 3.2$;
$4.9\times 3.2 = 16$।
\textbf{24.} प्रेरक वल्कुट से वह आदेश, और पेशियों के ग्राहियों से
वे संकेत, मेडुला में \omterm{def:b2:blood-pressure:baroreflex}{दाबग्राही} निवेश का भार फिर से बाँट देते
हैं ताकि वह प्रतिवर्त 95 के बजाय \qty{110}{mmHg} की रक्षा करे ---
यानी नियत बिंदु चढ़ाया गया है, वह पाश निष्क्रिय नहीं किया गया।
\textbf{25.} खड़े होने पर बचा हुआ पात \qty{5.7}{mmHg}; रक्तस्राव के बाद
\qty{86}{mmHg}; और व्यायाम में \qty{25}{L/min} तथा प्रति मिनट
\qty{4}{L} ऑक्सीजन।
\end{solution}
